\documentclass[11pt, a4paper]{article}

% Paquetes de idioma y tipografía
\usepackage[utf8]{inputenc}
\usepackage[spanish, es-tabla]{babel}
\usepackage[T1]{fontenc}
\usepackage{mathpazo} % Tipografía Palatino
\usepackage{microtype}

% Geometría de la página
\usepackage[margin=2.5cm, headheight=15pt]{geometry}

% Paquetes matemáticos
\usepackage{amsmath, amssymb, amsfonts, bm}

% Gráficos, colores y tablas
\usepackage{graphicx}
\usepackage[table, xcdraw]{xcolor}
\usepackage{booktabs}
\usepackage{tabularx}
\usepackage[most]{tcolorbox}
\usepackage{tikz}
\usetikzlibrary{shapes.geometric, arrows.meta, positioning, calc, babel}

% Personalización de listas y código
\usepackage{enumitem}
\usepackage{listings}

% Encabezados y pies de página
\usepackage{fancyhdr}
\pagestyle{fancy}
\fancyhf{}
\lhead{\textbf{UCB Sede Tarija} | Ing. Mecatrónica}
\chead{Robótica (IMT-342)}
\rhead{Clase 2 - Sem 3}
\lfoot{Prof: M.Sc. Ing. Bernardo Quiroga Turdera}
\rfoot{Página \thepage}
\renewcommand{\headrulewidth}{0.4pt}
\renewcommand{\footrulewidth}{0.4pt}

% Definición de colores institucionales
\definecolor{ucbblue}{RGB}{0, 51, 102}
\definecolor{ucbyellow}{RGB}{255, 184, 28}
\definecolor{codegray}{rgb}{0.95,0.95,0.95}

% Estilo para código Python
\lstdefinestyle{pythonstyle}{
    backgroundcolor=\color{codegray},
    commentstyle=\color{green!60!black},
    keywordstyle=\color{blue}\bfseries,
    numberstyle=\tiny\color{gray},
    stringstyle=\color{purple},
    basicstyle=\ttfamily\footnotesize,
    breakatwhitespace=false,         
    breaklines=true,                 
    captionpos=b,                    
    keepspaces=true,                 
    numbers=left,                    
    numbersep=5pt,                  
    showspaces=false,                
    showstringspaces=false,
    showtabs=false,                  
    tabsize=4,
    language=Python,
    frame=single,
    rulecolor=\color{gray!40}
}
\lstset{style=pythonstyle}

\begin{document}

\begin{center}
    \Large\textbf{\textcolor{ucbblue}{UNIVERSIDAD CATÓLICA BOLIVIANA "SAN PABLO" — SEDE TARIJA}}\\
    \vspace{0.2cm}
    \normalsize Departamento de Ciencias Exactas e Ingenierías | Carrera de Ingeniería Mecatrónica\\
    Asignatura: Robótica Industrial (IMT-342) \\
    Docente: M.Sc. Ing. Bernardo Quiroga Turdera
\end{center}

\vspace{0.4cm}
\begin{tcolorbox}[colback=ucbblue!5!white, colframe=ucbblue, title=\textbf{Documento de Cátedra: Semana 3 — Sesión 2}]
    \centering \large \textbf{Módulo Puente: De LERP a SLERP y \\ Guía de Estudio de Cuaterniones Unitarios}
\end{tcolorbox}
\vspace{0.4cm}

% =======================================================
% PARTE I: MÓDULO PUENTE LERP -> SLERP
% =======================================================
\section*{PARTE I: Módulo Puente — De LERP a SLERP (Fundamentación Paso a Paso)}

Si no se conoce LERP (\textit{Linear Interpolation}), intentar explicar SLERP directamente es una receta para la sobrecarga cognitiva. Pedagógicamente, SLERP no es más que la generalización geométrica de LERP sobre una superficie esférica ($\mathbb{S}^3$). 

\vspace{0.3cm}
\begin{center}
\begin{tikzpicture}[
    node distance=0.8cm and 0cm,
    box/.style={rectangle, draw=ucbblue, thick, fill=white, text width=12cm, rounded corners, align=left, minimum height=1cm, inner sep=10pt},
    arrow/.style={-Stealth, thick, ucbblue}
]
    \node[box, fill=gray!10, text centered, font=\bfseries] (title) {LA ESCALERA DIDÁCTICA DE LA INTERPOLACIÓN};
    
    \node[box, below=of title] (lerp) {\textbf{1. LERP Clásico (1D y $\mathbb{R}^3$):} $\quad p(t) = (1 - t)p_0 + t p_1$ \\ \small (Interpolación recta de posición: válida para traslación X, Y, Z)};
    
    \node[box, below=of lerp] (fail) {\textbf{2. El fallo de LERP en Rotación:} $\quad q_{\text{LERP}}$ corta la hiperesfera por dentro \\ \small ($\left\|q\right\| < 1$ deforma la orientación y el tamaño)};
    
    \node[box, below=of fail] (nlerp) {\textbf{3. NLERP (Normalized LERP):} $\quad q_{\text{NLERP}} = q_{\text{LERP}} / \left\|q_{\text{LERP}}\right\|$ \\ \small (Recupera la esfera, pero la velocidad angular NO es constante)};
    
    \node[box, below=of nlerp] (slerp) {\textbf{4. SLERP (Spherical LERP):} \\ \small (Recorre la geodésica con velocidad angular $\boldsymbol{\omega}$ constante)};

    \draw[arrow] (lerp) -- (fail);
    \draw[arrow] (fail) -- (nlerp);
    \draw[arrow] (nlerp) -- (slerp);
\end{tikzpicture}
\end{center}

\subsection*{1. ¿Qué es LERP? (Explicación Intuitiva para Pizarra)}

\textbf{A. Intuición en 1D: La Palanca o Porcentaje de Avance}\\
Estás en la posición $p_0 = 10\text{ cm}$ a las $t = 0$ y quieres llegar a $p_1 = 50\text{ cm}$ a las $t = 1$. ¿Dónde estás a la mitad del camino ($t = 0.5$)? Intuitivamente: $30\text{ cm}$.

La ecuación de la recta que conecta dos valores con un parámetro de tiempo normalizado $t \in [0, 1]$ es:
\begin{equation*}
    p(t) = p_0 + t(p_1 - p_0)
\end{equation*}
Reordenando los términos para ver los pesos de ponderación:
\begin{equation}
    \mathbf{p(t) = (1 - t)p_0 + t p_1}
\end{equation}
\begin{itemize}[label=$\circ$]
    \item Cuando $t = 0$: $p(0) = 1(p_0) + 0(p_1) = p_0$ (Punto inicial).
    \item Cuando $t = 1$: $p(1) = 0(p_0) + 1(p_1) = p_1$ (Punto final).
    \item Cuando $t = 0.5$: $p(0.5) = 0.5(p_0) + 0.5(p_1)$ (Promedio aritmético exacto).
\end{itemize}

\begin{tcolorbox}[colback=yellow!10!white, colframe=ucbyellow!80!black, title=Concepto Clave para el Alumno]
LERP es una combinación convexa: los pesos $(1-t)$ y $t$ siempre suman $1$. Es como una perilla de volumen (fader) que desvanece suavemente el estado inicial y amplifica el estado final.
\end{tcolorbox}

\textbf{B. LERP en $\mathbb{R}^3$: Trayectoria Rectilínea del Extremo del Robot}\\
En robótica cartesiana, para mover la pinza neumática del robot desde un punto $A = [x_0, y_0, z_0]^T$ hasta un punto $B = [x_1, y_1, z_1]^T$ en línea recta, aplicamos LERP vectorialmente:
\begin{equation*}
    \mathbf{p}(t) = (1 - t)\mathbf{p}_0 + t\mathbf{p}_1 \quad \text{con } t \in [0, 1]
\end{equation*}
Esto traza el segmento de recta euclidiano que une ambos puntos en el espacio tridimensional de trabajo.

\subsection*{2. El Experimento Mental: ¿Qué pasa si aplicamos LERP a la Orientación?}
Pídeles que imaginen dos orientaciones representadas por cuaterniones unitarios $\mathbf{q}_0$ y $\mathbf{q}_1$ (ambos de longitud $\left\|\mathbf{q}\right\| = 1$, viviendo en la superficie de la hiperesfera $\mathbb{S}^3$).
Si intentamos hacer LERP directo sobre los cuaterniones: $\mathbf{q}_{\text{LERP}}(t) = (1 - t)\mathbf{q}_0 + t\mathbf{q}_1$

\vspace{0.4cm}
\begin{center}
\begin{tikzpicture}[scale=1.5, thick]
    % Esfera S^3
    \draw[fill=blue!5] (0,0) circle (2cm);
    \draw[dashed, gray] (0,0) circle (2cm);
    \node at (-1.5, 1.5) {$\mathbb{S}^3$};
    
    % Vectores q0 y q1
    \coordinate (O) at (0,0);
    \coordinate (q0) at (0:2cm);
    \coordinate (q1) at (90:2cm);
    
    \draw[-Stealth, ucbblue] (O) -- (q0) node[right] {$\mathbf{q}_0$ ($\left\|\mathbf{q}\right\|=1$)};
    \draw[-Stealth, ucbblue] (O) -- (q1) node[above] {$\mathbf{q}_1$ ($\left\|\mathbf{q}\right\|=1$)};
    
    % Línea LERP (Cuerda)
    \draw[red!80!black, thick] (q0) -- (q1);
    \coordinate (lerp_mid) at (1,1);
    \filldraw[red] (lerp_mid) circle (1.5pt) node[right, xshift=2pt] {$\mathbf{q}_{\text{LERP}}(0.5)$ (Corta por dentro, $\left\|\mathbf{q}\right\| = 0.707$)};
    
    % Proyección NLERP
    \draw[-Stealth, dashed, green!60!black] (O) -- (lerp_mid);
    \draw[-Stealth, green!60!black, thick] (lerp_mid) -- (45:2cm);
    \filldraw[green!50!black] (45:2cm) circle (1.5pt) node[above right] {$\mathbf{q}_{\text{NLERP}}(0.5)$};
    
    % Arco SLERP
    \draw[ucbblue, very thick, -Stealth] (q0) arc (0:90:2cm);
    \node[ucbblue, fill=blue!5] at (-0.5, -0.5) {SLERP (Superficie)};
\end{tikzpicture}
\end{center}

\textbf{¿Cuáles son las dos fallas físicas fatales de este enfoque?}\\
\textbf{Falla 1: Colapso de la Norma Unitaria ($\left\|\mathbf{q}\right\| < 1$).} La línea recta que une $\mathbf{q}_0$ y $\mathbf{q}_1$ corta la esfera por dentro. En el punto medio ($t = 0.5$), la longitud del cuaternión cae. Si intentas convertir $\mathbf{q}_{\text{LERP}}$ a matriz de rotación, la matriz deja de ser ortogonal y deforma/achica el robot en la simulación.

\textbf{Falla 2 (NLERP): Pérdida de la Velocidad Angular Constante.} Para corregir la norma, se divide entre su magnitud (NLERP):
\begin{equation*}
    \mathbf{q}_{\text{NLERP}}(t) = \frac{(1 - t)\mathbf{q}_0 + t\mathbf{q}_1}{\left\| (1 - t)\mathbf{q}_0 + t\mathbf{q}_1 \right\|}
\end{equation*}
El problema: Aunque ahora tiene norma $1$, la velocidad angular no es uniforme. El robot gira lento cerca de los extremos y se acelera bruscamente en el centro, generando \textit{jerk} y vibraciones mecánicas en la muñeca del ABB IRB 120.

\subsection*{3. La Necesidad de SLERP: El Gran Círculo}
"Para que la velocidad de rotación sea rigurosamente constante ($\left\|\boldsymbol{\omega}\right\| = \text{cte}$), no podemos viajar en línea recta y proyectar. Debemos viajar directamente a lo largo del arco esférico. SLERP es la fórmula trigonométrica que calcula ese viaje sobre el círculo máximo."
\begin{equation}
    \text{SLERP}(\mathbf{q}_0, \mathbf{q}_1; t) = \frac{\sin((1 - t)\Omega)}{\sin\Omega}\mathbf{q}_0 + \frac{\sin(t\Omega)}{\sin\Omega}\mathbf{q}_1
\end{equation}

La estructura es idéntica a LERP, pero los pesos $\frac{\sin((1-t)\Omega)}{\sin\Omega}$ y $\frac{\sin(t\Omega)}{\sin\Omega}$ compensan la curvatura esférica para que el avance angular sea perfectamente lineal en el tiempo.

\subsection*{4. Cuadro Comparativo (Para Diapositiva)}
\begin{table}[h!]
    \centering
    \renewcommand{\arraystretch}{1.5}
    \resizebox{\textwidth}{!}{
    \begin{tabular}{l p{4cm} p{2.5cm} p{3.5cm} p{4cm}}
    \toprule
    \textbf{Método} & \textbf{Fórmula Matemática} & \textbf{¿Permanece en $\mathbb{S}^3$?} & \textbf{¿Velocidad $\boldsymbol{\omega}$ Constante?} & \textbf{Uso Típico en Ingeniería} \\ \midrule
    \textbf{LERP} & $\mathbf{p}(t) = (1-t)\mathbf{p}_0 + t\mathbf{p}_1$ & No (Corta por dentro) & N/A (No aplica a rotación) & Posición XYZ de la punta del robot. \\ 
    \textbf{NLERP} & $\mathbf{q}(t) = \text{LERP} / \left\|\text{LERP}\right\|$ & Sí (Normalizado) & No (Acelera y frena) & Gráficos en videojuegos. \\ 
    \textbf{SLERP} & $\frac{\sin((1-t)\Omega)}{\sin\Omega}\mathbf{q}_0 + \frac{\sin(t\Omega)}{\sin\Omega}\mathbf{q}_1$ & Sí (Geodésica pura) & Sí ($\left\|\boldsymbol{\omega}\right\| = \text{constante}$) & Robótica Industrial, ROS. \\ \bottomrule
    \end{tabular}
    }
\end{table}

\subsection*{5. Script Demostrativo Corto en Python (Para proyectar en 2 minutos)}
\begin{lstlisting}[language=Python]
import numpy as np

# Dos orientaciones separadas 90 grados alrededor de Z
q0 = np.array([1.0, 0.0, 0.0, 0.0])              # Rotacion 0 grados
q1 = np.array([np.cos(np.pi/4), 0, 0, np.sin(np.pi/4)]) # Rotacion 90 grados
t = 0.5  # Punto medio

# 1. LERP
q_lerp = (1 - t)*q0 + t*q1
print(f"Norma LERP a t=0.5 : {np.linalg.norm(q_lerp):.4f} (!Pierde norma unitaria!)")

# 2. NLERP
q_nlerp = q_lerp / np.linalg.norm(q_lerp)
print(f"Norma NLERP a t=0.5: {np.linalg.norm(q_nlerp):.4f} (Recupera norma, omega variable)")

# 3. SLERP (Shoemake)
omega = np.arccos(np.dot(q0, q1))
q_slerp = (np.sin((1-t)*omega)/np.sin(omega))*q0 + (np.sin(t*omega)/np.sin(omega))*q1
print(f"Norma SLERP a t=0.5: {np.linalg.norm(q_slerp):.4f} (Norma perfecta y omega constante)")
\end{lstlisting}

\newpage
% =======================================================
% PARTE II: GUÍA DE ESTUDIO Y FUNDAMENTACIÓN MATEMÁTICA
% =======================================================
\section*{PARTE II: Guía de Estudio y Fundamentación Matemática}
\textit{Cuaterniones Unitarios, Operador Sándwich y SLERP}

\subsection*{Introducción y Propósito Didáctico}
El propósito de este documento es desmitificar y justificar geométricamente los tres conceptos más densos de la orientación espacial:
\begin{enumerate}
    \item \textbf{El Operador Sándwich ($\mathbf{q} \otimes \tilde{p} \otimes \mathbf{q}^*$):} Por qué una rotación requiere multiplicar dos veces en 4D y cómo esto reconstruye la fórmula física de Euler-Rodrigues.
    \item \textbf{La Matriz $R(\mathbf{q}) \in SO(3)$ y la Doble Cobertura ($\mathbf{q} \equiv -\mathbf{q}$):} De dónde surgen sus 9 componentes y el riesgo industrial de ignorar el signo antipodal.
    \item \textbf{Interpolación Esférica Lineal (SLERP):} La matemática de las trayectorias con velocidad angular uniforme.
\end{enumerate}

\subsection*{SECCIÓN 1: El Operador Sándwich de Rotación Activa}

\begin{center}
\begin{tikzpicture}[node distance=2.2cm, auto, thick, -Stealth]
    \node[draw=ucbblue, rounded corners, fill=white, text width=6cm, align=center, inner sep=5pt] (v1) {\textbf{[Vector Físico 3D: $p$]}};
    
    \node[draw=ucbblue, rounded corners, fill=gray!10, text width=6cm, align=center, inner sep=5pt, below=0.7cm of v1] (v2) {\textbf{[Cuaternión Puro: $\tilde{p} = [0, p] \in \mathbb{H}$]}\\ \small (Promoción: se añade parte real $= 0$)};
    
    \node[draw=red!80!black, rounded corners, fill=red!10, text width=7.5cm, align=center, inner sep=5pt, below=0.7cm of v2] (v3) {\textbf{[Cuaternión Intermedio: $\mathbf{q} \otimes \tilde{p}$]} \\ \small ¡Parte Escalar $\neq 0$! (Escapa a 4D)};
    
    \node[draw=green!60!black, rounded corners, fill=green!10, text width=7.5cm, align=center, inner sep=5pt, below=0.7cm of v3] (v4) {\textbf{[Cuaternión Resultante: $\tilde{p}' = [0, p'] \in \mathbb{H}$]} \\ \small Multiplicación por $\mathbf{q}^*$. ¡Parte Escalar $= 0$!};
    
    \node[draw=ucbblue, rounded corners, fill=white, text width=6cm, align=center, inner sep=5pt, below=0.7cm of v4] (v5) {\textbf{[Vector Rotado Físico 3D: $p'$]}\\ \small (Extracción Vectorial)};

    \draw (v1) -- (v2);
    \draw (v2) -- (v3);
    \draw (v3) -- (v4);
    \draw (v4) -- (v5);
\end{tikzpicture}
\end{center}

\subsubsection*{1.1 La Intuición Geométrica}
Cualquier orientación de un cuerpo en $\mathbb{R}^3$ equivale a un giro $\theta$ sobre un eje unitario $\hat{u}$. Un vector $p$ se descompone en componente paralela $p_\parallel = (\hat{u} \cdot p)\hat{u}$ y perpendicular $p_\perp = p - (\hat{u} \cdot p)\hat{u}$.
En 1840, Olinde Rodrigues demostró que el vector resultante rotado $p'$ es:
\begin{equation}
    \mathbf{p}' = \mathbf{p} \cos\theta + (\hat{\mathbf{u}} \times \mathbf{p})\sin\theta + \hat{\mathbf{u}}(\hat{\mathbf{u}} \cdot \mathbf{p})(1 - \cos \theta)
\end{equation}


\subsubsection*{1.2 El Dilema de la Cuarta Dimensión ($\mathbf{q} \otimes \tilde{p}$)}
Un cuaternión unitario de rotación es $\mathbf{q} = \left[ \cos(\frac{\theta}{2}), \; \hat{u} \sin(\frac{\theta}{2}) \right] = [q_w, \; \mathbf{v}]$. Si rotamos multiplicando solo por la izquierda:
\begin{equation*}
    \mathbf{q} \otimes \tilde{p} = [q_w, \; \mathbf{v}] \otimes [0, \; p] = \begin{bmatrix} q_w (0) - \mathbf{v} \cdot p \\ q_w p + 0\mathbf{v} + \mathbf{v} \times p \end{bmatrix} = \begin{bmatrix} \mathbf{-\mathbf{v} \cdot p} \\ q_w p + \mathbf{v} \times p \end{bmatrix}
\end{equation*}
\textbf{El problema físico:} La parte escalar es $-\mathbf{v} \cdot p \neq 0$. El vector "escapó" de nuestro espacio físico 3D a 4D.

\subsubsection*{1.3 La Acción del Conjugado ($\mathbf{q}^*$)}
Para corregir la fuga, Hamilton post-multiplicó por el conjugado $\mathbf{q}^* = [q_w, \; -\mathbf{v}]$:
\begin{equation*}
    \tilde{p}' = (\mathbf{q} \otimes \tilde{p}) \otimes \mathbf{q}^* = \begin{bmatrix} -\mathbf{v} \cdot p \\ q_w p + \mathbf{v} \times p \end{bmatrix} \otimes \begin{bmatrix} q_w \\ -\mathbf{v} \end{bmatrix}
\end{equation*}
\textbf{Parte Escalar de $\tilde{p}'$:}
\begin{align*}
    \text{Escalar} &= (-\mathbf{v} \cdot p)(q_w) - \left( q_w p + \mathbf{v} \times p \right) \cdot (-\mathbf{v}) \\
    &= -q_w(\mathbf{v} \cdot p) + q_w(p \cdot \mathbf{v}) + \underbrace{(\mathbf{v} \times p) \cdot \mathbf{v}}_{= 0 \text{ (ortogonales)}} = \mathbf{0}
\end{align*}
\textbf{Parte Vectorial ($p'$):}
\begin{equation*}
    p' = (-\mathbf{v} \cdot p)(-\mathbf{v}) + q_w(q_w p + \mathbf{v} \times p) + (q_w p + \mathbf{v} \times p) \times (-\mathbf{v})
\end{equation*}
Expandiendo y usando la identidad del triple producto vectorial $(A \times B) \times C = (A \cdot C)B - (B \cdot C)A$:
\begin{equation}
    \mathbf{p' = (q_w^2 - \left\|\mathbf{v}\right\|^2)p + 2(\mathbf{v} \cdot p)\mathbf{v} + 2q_w(\mathbf{v} \times p)}
\end{equation}

\subsubsection*{1.4 Equivalencia con Rodrigues}
Sustituyendo $q_w = \cos(\theta/2)$ y $\mathbf{v} = \hat{u}\sin(\theta/2)$:
\begin{itemize}
    \item $q_w^2 - \left\|\mathbf{v}\right\|^2 = \cos^2(\theta/2) - \sin^2(\theta/2) = \cos\theta$
    \item $2q_w(\mathbf{v} \times p) = (\hat{u} \times p)\sin\theta$
    \item $2(\mathbf{v} \cdot p)\mathbf{v} = 2\sin^2(\theta/2)(\hat{u} \cdot p)\hat{u} = (1 - \cos\theta)(\hat{u} \cdot p)\hat{u}$
\end{itemize}
Uniendo los términos, recuperamos exactamente la Fórmula de Rotación de Rodrigues.

\subsection*{SECCIÓN 2: Conversión a $SO(3)$ y Fenómeno de Doble Cobertura}

\subsubsection*{2.1 Deducción de la Matriz $R(\mathbf{q}) \in SO(3)$}
Evaluando el sándwich en $\hat{e}_1 = [1, 0, 0]^T$ obtenemos la primera columna de la matriz:
\begin{equation*}
    R_{\cdot 1} = \begin{bmatrix} q_w^2 + q_x^2 - q_y^2 - q_z^2 \\ 2(q_x q_y + q_w q_z) \\ 2(q_x q_z - q_w q_y) \end{bmatrix} = \begin{bmatrix} \mathbf{1 - 2(q_y^2 + q_z^2)} \\ 2(q_x q_y + q_w q_z) \\ 2(q_x q_z - q_w q_y) \end{bmatrix}
\end{equation*}
Repitiendo para $\hat{e}_2$ y $\hat{e}_3$, formamos los 9 términos de $R(\mathbf{q})$.

\subsubsection*{2.2 Doble Cobertura ($\mathbf{q} \equiv -\mathbf{q}$)}
Como todos los 9 términos de $R(\mathbf{q})$ son cuadráticos o productos cruzados, al evaluar $-\mathbf{q}$:
\begin{equation}
    \mathbf{R(-\mathbf{q}) \equiv R(\mathbf{q})}
\end{equation}

\begin{center}
\begin{tikzpicture}[node distance=3cm, auto, thick]
    \node[draw=ucbblue, rounded corners, fill=blue!5, text width=4cm, align=center, minimum height=1.5cm] (s3) {\textbf{Hiperesfera $\mathbb{S}^3$ (4D)}\\ Polo Norte: $+\mathbf{q}$ \\ Polo Sur: $-\mathbf{q}$};
    \node[draw=black, rounded corners, fill=gray!10, text width=4cm, align=center, minimum height=1.5cm, right=of s3] (so3) {\textbf{Espacio $SO(3)$ (3D)}\\ Orientación Física Única};
    
    \draw[-Stealth, ucbblue] (s3.15) -- node[above] {Proyectan a} (so3.165);
    \draw[-Stealth, ucbblue] (s3.345) -- node[below] {misma matriz} (so3.195);
    
    \node[below=0.2cm of s3] {(\textit{Doble Cobertura 2:1})};
\end{tikzpicture}
\end{center}

\subsubsection*{2.3 El Peligro de Ingeniería en el ABB IRB 120}
Si el robot está en $\mathbf{q}_0 = [1.0, 0, 0, 0]$ y la meta es $\mathbf{q}_1 = [-0.9999, 0, 0, -0.0100]$ (desviación de $\approx 1.15^\circ$). Físicamente es un giro mínimo. Pero si el software no verifica el producto punto $\mathbf{q}_0 \cdot \mathbf{q}_1 = -0.9999 < 0$, interpretará que debe rotar $360^\circ - 1.15^\circ = 358.85^\circ$, causando daños mecánicos y disparos de sobrecorriente.

\subsection*{SECCIÓN 3: SLERP — Deducción Geométrica}
\subsubsection*{3.1 Analogía Intuitiva}
LERP vuela en un ``túnel'' subterráneo. SLERP vuela sobre la superficie terrestre siguiendo un Gran Círculo (Geodésica), manteniendo el radio constante ($\left\|\mathbf{q}\right\| = 1$).

\vspace{0.3cm}
\begin{center}
\begin{tikzpicture}[scale=1.5, thick]
    % Sector circular
    \draw[fill=blue!5] (0,0) -- (0:3cm) arc (0:60:3cm) -- cycle;
    
    \coordinate (O) at (0,0);
    \coordinate (q0) at (0:3cm);
    \coordinate (q1) at (60:3cm);
    \coordinate (qt) at (25:3cm);
    
    \draw[-Stealth, ucbblue, very thick] (O) -- (q0) node[right] {$\mathbf{q}_0$};
    \draw[-Stealth, ucbblue, very thick] (O) -- (q1) node[above right] {$\mathbf{q}_1$};
    \draw[-Stealth, red!80!black, very thick] (O) -- (qt) node[right, xshift=2pt] {$\mathbf{q}(t)$};
    
    % Ángulos
    \draw[dashed, gray] (0:1cm) arc (0:60:1cm);
    \node at (15:1.2cm) {$t\Omega$};
    \node at (45:1.2cm) {$(1-t)\Omega$};
\end{tikzpicture}
\end{center}

\subsubsection*{3.2 Deducción Geométrica de la Ecuación}
Expresamos $\mathbf{q}(t)$ como combinación lineal en el plano: $\mathbf{q}(t) = a(t)\mathbf{q}_0 + b(t)\mathbf{q}_1$.
Aplicando la Ley de los Senos en el triángulo del paralelogramo formado:
\begin{equation}
    \text{SLERP}(\mathbf{q}_0, \mathbf{q}_1; t) = \frac{\sin((1 - t)\Omega)}{\sin\Omega}\mathbf{q}_0 + \frac{\sin(t\Omega)}{\sin\Omega}\mathbf{q}_1
\end{equation}


\subsubsection*{3.3 Algoritmo Completo (Shortest Path)}
\begin{center}
\begin{tikzpicture}[
    node distance=1cm and 1cm,
    box/.style={rectangle, draw=black, thick, fill=white, text width=4cm, align=center, rounded corners, minimum height=0.8cm},
    decision/.style={diamond, draw=ucbblue, thick, fill=blue!5, text width=2cm, align=center, inner sep=0pt},
    arrow/.style={-Stealth, thick}
]
    \node[box] (input) {Entrada: $\mathbf{q}_0, \mathbf{q}_1$, $t$};
    \node[box, below=0.6cm of input] (dot) {dot = $\mathbf{q}_0 \cdot \mathbf{q}_1$};
    \node[decision, below=0.6cm of dot] (check1) {dot < 0 ?};
    
    \node[box, right=1.5cm of check1] (inv) {$\mathbf{q}_1 = -\mathbf{q}_1$ \\ dot = -dot};
    \node[decision, below=0.8cm of check1] (check2) {dot > 0.999 ?};
    
    \node[box, left=1cm of check2, yshift=-1.5cm] (lerp) {LERP: \\ $(1-t)\mathbf{q}_0 + t\mathbf{q}_1$ \\ (Normalizar)};
    \node[box, right=1cm of check2, yshift=-1.5cm] (slerp) {SLERP: \\ $\Omega = \arccos(\text{dot})$ \\ Eq. Trigonométrica};
    
    \node[box, fill=green!10, below=3.5cm of check2] (output) {Retornar $\mathbf{q}(t) \in \mathbb{S}^3$};

    \draw[arrow] (input) -- (dot);
    \draw[arrow] (dot) -- (check1);
    \draw[arrow] (check1) -- node[above]{Sí} (inv);
    \draw[arrow] (inv) |- (check2);
    \draw[arrow] (check1) -- node[left]{No} (check2);
    
    \draw[arrow] (check2) -| node[above, near start]{Sí ($\Omega \approx 0$)} (lerp);
    \draw[arrow] (check2) -| node[above, near start]{No} (slerp);
    
    \draw[arrow] (lerp) |- (output);
    \draw[arrow] (slerp) |- (output);
\end{tikzpicture}
\end{center}

\subsection*{SECCIÓN 4: Preguntas de Autoevaluación Conceptual}
\begin{tcolorbox}[colback=white, colframe=ucbblue, title=Verificación de Conocimientos]
\begin{enumerate}
    \item ¿Por qué al rotar un vector $p \in \mathbb{R}^3$ con cuaterniones se debe multiplicar por la izquierda por $\mathbf{q}$ y por la derecha por $\mathbf{q}^*$, en lugar de multiplicar una sola vez?
    \item Si un sensor de visión publica el cuaternión $\mathbf{q}_A = [0, 0, 1, 0]$ y otro algoritmo publica $\mathbf{q}_B = [0, 0, -1, 0]$, ¿el robot adoptará posturas físicas diferentes al alcanzar ambas metas? Justifique analíticamente su respuesta.
    \item En el algoritmo SLERP, ¿qué problema físico ocurriría en el manipulador ABB IRB 120 si omitimos la condición \texttt{if dot < 0: dot = -dot; q1 = -q1}?
    \item ¿Por qué no es recomendable interpolar linealmente los 3 ángulos de Euler $(\phi, \theta, \psi)$ entre dos puntos de paso de una trayectoria robótica? Mencione al menos dos consecuencias físicas sobre los actuadores.
\end{enumerate}
\end{tcolorbox}

\end{document}
